\documentclass{article}
\usepackage{geometry}
\usepackage{graphicx} % Required for inserting images
\usepackage{amsmath, amsthm, amssymb}
\usepackage{parskip}
\newgeometry{vmargin={15mm}, hmargin={24mm,34mm}}
\theoremstyle{definition} 
\newtheorem{definition}{Definition}

\newtheorem{theorem}{Theorem}[section]
\newtheorem{lemma}[theorem]{Lemma}
\newtheorem{proposition}[theorem]{Proposition}
\newtheorem{example}[theorem]{Example}
\newtheorem{corollary}{Corollary}[theorem]
\newcommand{\nil}{\operatorname{nil}}

\title{Localization}
\date{April 2025}
\author{Boran Erol}

\begin{document}

\maketitle

The following is ATOCA Chapter 11, Exercise 5.

\begin{lemma}
    Let $ R^{\prime },R^{\prime \prime} $ be rings.
    Let $ R := R^{\prime} \times R^{\prime \prime} $.
    Let $ S = \{ (1,1), (1,0) \} $.
    Then, $ R^{\prime} $ is the localization of $ R $ at $ S $.
\end{lemma}
\begin{proof}
    We have a ring homomorphism from $ R $ onto $ R^{\prime} $
    given by the projection map with kernel $ R^{\prime \prime} $.

    By ATOCA 11.10, $ \ker(\phi_{S}) = \langle 0 \rangle^{S} = R^{\prime \prime} $.

    Then, the induced map from $ S^{-1}R $ to $ R^{\prime} $ is injective.
    The map is also surjective since the diagram given in ATOCA 11.3 commutes.
    Thus, the two rings are isomorphic.
\end{proof}

Let $ R $ be a ring and $ S $ be a multiplicative subset of $ R $.
Let $ J $ be an ideal of $ S^{-1}R $.
Let $ \phi = \phi_{S} $.
Let $ I = \phi^{-1}(J) $.

\begin{lemma}
    $ I $ is saturated.
\end{lemma}
\begin{proof}
    Let $ x \in I^{S} $. Then, there exists $ s \in S $ such that $ xs \in I $. 
    Then, $ \phi(xs) \in J $. But $ \phi(s) $ is a unit in $ S^{-1}R $, and thus
    we have that $ \phi(x) = \phi(x) \phi(s) \phi(s^{-1}) = \phi(xs) \phi(s^{-1}) \in J $.
    Then, $ x \in I $.
\end{proof}

\begin{lemma}
    $ J = I (S^{-1}R) $.
\end{lemma}
\begin{proof}
    Let $ \frac{x}{s} \in J $. Then, $ \frac{x}{1} \in J $.
    Then, $ x \in I $. Then, $ \frac{x}{1} \frac{1}{s} \in I (S^{-1}R) $.
    Thus, $ J \subseteq I (S^{-1}R) $.

    Let $ x \in I $. Then, $ \frac{x}{1} \in J $.
    But then $ J $ contains $ I (S^{-1}R) $ as $ J $ is closed under multiplication by $ S^{-1}R $.
\end{proof}

Now, let $ I $ be an ideal of $ R $.

\begin{lemma}
    $ I (S^{-1}R) $ is a saturated ideal of $ S^{-1}R $.
\end{lemma}
\begin{proof}
    Let $ x \in (I (S^{-1}R))^{S} $.
    Then, there exists some $ s \in S $ such that $ xs \in I (S^{-1}R) $.
    Since $ s $ is a unit in $ S^{-1}R $, this implies that $ x \in I (S^{-1}R) $.
\end{proof}

\begin{lemma}
    The preimage of $ I (S^{-1}R) $ under $ \phi $ is the
    saturation of $ I $.
\end{lemma}
\begin{proof}
    We show both inclusions.

    Let $ x \in I^{S} $. Then, there exists some $ s \in S $
    such that $ xs \in I $. Then, $ \phi(xs) \in J $, so $ \phi(x) \in J $.
    Since $ J = I (S^{-1}R) $, we have the desired result.

    Let $ x $ be in the preimage of $ I (S^{-1}R) $.
    Then, $ \phi(x) \in I (S^{-1}R) $, so $ \frac{x}{1} = \frac{y}{s} $
    for some $ y \in I $ and $ s \in S $. Then, there is some $ t \in S $
    such that $ xst = yt \in I $ since $ y \in I $.
    Then, $ s \in I^{s} $.
\end{proof}

Let $ R $ be a ring and $ S $ be a multiplicative subset of $ R $.
Let $ P $ be a prime ideal of $ R $.

Let $ x \in P^{S} $. Then, there exists some $ s \in S $ such that $ xs \in P $.
Then, either $ x \in P $ or $ s \in P $. Since $ s \notin P $, $ x \in P $.
Thus, $ P = P^{S} $.

\begin{lemma}
    $ P (S^{-1}R) $ is a prime ideal.
\end{lemma}
\begin{proof}
    Since the preimage of $ P (S^{-1}R) $ is $ P^{S} = P $, $ P (S^{-1}R) $ can't be a unit ideal.
    
    Now, let $ \frac{a}{s} \cdot \frac{b}{t} \in P (S^{-1}R) $.
    Then, $ \frac{ab}{st} \in P (S^{-1}R) $ and thus $ ab \in P $,
    and thus $ a \in P $ or $ b \in P $, so $ \frac{a}{s} \in P (S^{-1}R) $
    or $ \frac{b}{t} \in P (S^{-1}R) $.
\end{proof}

We thus gave a bijective correspondece that respects inclusions
between saturated ideals of $ R $ and ideals of $ S^{-1}R $.
Notice that this bijection also respects prime ideals, which are
always saturated.

The following is Exercise 11.27(c) from ATOCA.

\begin{lemma}
    Let $ R $ be a ring, $ S $ be a multiplicative subset of $ R $.
    Let $ I,J $ be ideals of $ R $.
    Then, $ (I^{S}J^{S})^{S} = (IJ)^{S} $.
\end{lemma}
\begin{proof}
    
\end{proof}

The following is from ATOCA 12.11 and beyond.

Let $ R $ be a ring, $ S $ be a multiplicative subset of $ R $ and $ N $ be a submodule
of $ M $.

Let's first show that $ N^{S} $ is a submodule of $ M $.

Let $ x,y \in N^{S} $. Then, there exists $ s,t \in S $ such
that $ sx + ty \in N $. Then, $ (x + y)st = xst + yst \in N $
so $ x + y \in N^{S} $. Moreover, for any $ r \in R $, $ rxs \in N $.
Thus, $ N^{S} $ is a submodule. This gives us ATOCA 12.12(1a).

Notice that 12.12(1b) is an immediate consequence of the exactness of localization.

\begin{lemma}[ATOCA 12.12(3b)]
    $ S^{-1} N = S^{-1} N^{S} $.
\end{lemma}
\begin{proof}
    Let $ \frac{x}{s} \in S^{-1} N^{S} $.

    Since $ x \in N^{S} $, there exists some $ t \in S $ such that $ tx \in N $.
    Then, $ \frac{tx}{s} \in N $ so $ \frac{x}{s} = \frac{1}{t} \frac{tx}{s} \in N $.
\end{proof}



\begin{lemma}[ATOCA 12.17]
    Let $ R $ be a ring and $ S $ be a multiplicative subset of $ R $.
    Let $ M $ be an $ R $-module.

    Then, $ S^{-1} Ann(M) \subset Ann(S^{-1} M) $.
\end{lemma}
\begin{proof}
    
\end{proof}

\begin{definition}
    Let $ R $ be a ring and $ M $ be an $ R $-module.
    $ x $ is \textbf{nilpotent on M} if there exists
    some $ n \geq 1 $ such that $ x^{n}M  = 0 $.
    In other words, $ x \in \sqrt{Ann(M)} $.

    We denote the set of such elements by $ \nil(M) $.
\end{definition}

Notice that if $ M = R $, we recover the usual definition of a nilpotent element.

\begin{lemma}[ATOCA 12.23]
    Let $ R $ be a ring and $ I $ be an ideal of $ I $. Then, $ \nil(R/I) = \sqrt{I} $.
\end{lemma}

\begin{proposition}
    Let $ R $ be a ring and $ S $ be a multiplicatively closed subset of $ R $.
    Let $ Q \subsetneq M $ be $ R $-modules.

    Assume $ S \cap \nil(M/Q) \neq \varnothing $.

    Then, $ Q^{S} = M $ and $ S^{-1}Q = S^{-1}M $.
\end{proposition}
\begin{proof}
    Let $ s \in S \cap \nil(M/Q) $.

    Then, there exists some $ n \geq 1 $ such that for every $ m \in M $,
    $ s^{n}m \in Q $. Since $ s^{n} \in S $, this implies that $ M = Q^{S} $.

    Since $ S^{-1}Q = S^{-1} Q^{S} $, we have the desired result.
\end{proof}

% TODO: I think it follows from 12.17.
\begin{lemma}[ATOCA 12.40]
    Let $ R $ be a ring and $ S $ be a multiplicative subset of $ R $.
    Let $ M $ be an $ R $-module.

    Then, $ S^{-1} (\nil(M)) \subseteq \nil(S^{-1}M) $.
\end{lemma}
\begin{proof}
    Let $ \frac{x}{s} \in S^{-1} (\nil(M)) $ where $ x \in \nil(M) $
    and $ s \in S $.

    Then, there exists some $ n \geq 1 $ such that $ x^{n} $

    Notice that $ x \in $
\end{proof}

\end{document}
